\ChapterImageStar[cap:Marco-Teorico]{Marco Teórico}{./images/fondo.png}\label{cap:marcoTeorico}
\mbox{}\\
%------------------------------------------------------------------------------------------------------------------------------------
\noindent
El marco teórico presenta los fundamentos investigativos y normativos que respaldan la propuesta, abordando estudios relacionados con
la adopción de la computación en la nube y la seguridad de la información.

\section{Adopción de cloud en universidades}
\noindent
La adopción de la computación en la nube en instituciones académicas ha incrementado significativamente en los últimos años, debido a su 
capacidad para mejorar la gestión de recursos tecnológicos y facilitar el acceso a servicios digitales. En este sentido,~\cite{aydin2021cloud}
señala que la computación en la nube no solo responde a necesidades académicas, sino que también actúa como una guía para la migración y 
administración de datos institucionales, contribuyendo a la continuidad operativa.

\section{Crecimiento y retos de la nube}
\noindent
El crecimiento de los servicios en la nube ha sido impulsado por su flexibilidad y escalabilidad. Sin embargo, este avance también ha traído consigo
nuevos retos en materia de seguridad. De acuerdo con~\cite{ibm2025cloud} la computación en la nube ofrece múltiples beneficios, la seguridad sigue 
siendo uno de los principales desafíos, ya que requiere la implementación de capacidades técnicas especializadas y procedimientos adecuados.

\section{Seguridad en la nube basada en riesgos}
\noindent
La seguridad en la nube debe abordarse desde un enfoque basado en riesgos. En este sentido, la \textit{European Network and Information Security Agency}
propone que la protección en entornos cloud deben fundamentarse en la identificación, análisis y mitigación de riesgos asociados a las amenazas y 
vulnerabilidades propias de estos sistemas. Esto implica que la adopción de estrategias de seguridad en la nube debe orientarse a la implementación 
de controles adaptativos y mecanismos de monitoreo continuo que permitan responder de manera efectiva a los riesgos identificados.

\section{Modelo de seguridad (CIA)}
\noindent
En el ámbito de la seguridad de la información, la tríada \CIA \textit{(Confidentiality, Integrity, Availability)} constituye uno de los modelos 
fundamentales para el diseño y evaluación de mecanismos de protección. Más allá de su definición conceptual, este modelo sirve como base para la 
identificación de requisitos de seguridad y la priorización de controles en sistemas computacionales. Según~\cite{stallings2012computer}, estos tres 
principios permiten estructurar las estrategias de seguridad en función de los riesgos a los que se encuentran expuestos los activos de información. 
En este sentido, la confidencialidad se relaciona con la restricción del acceso a usuarios autorizados, la integridad con la protección frente a 
modificaciones no autorizadas y la disponibilidad con la garantía de acceso oportuno a la información.

En entornos de computación en la nube, la aplicación de la tríada CIA adquiere una mayor complejidad debido a factores como la distribución de los 
recursos, la multi-tenencia y la externalización de servicios. Esto implica que los mecanismos de seguridad deben adaptarse para garantizar estos principios 
en escenarios donde el control de la infraestructura es compartido o delegado a terceros.

Por tanto, la tríada CIA no solo actúa como un marco conceptual, sino también como un criterio práctico para el diseño de arquitecturas de seguridad, 
permitiendo evaluar la efectividad de los controles implementados y su alineación con los objetivos de protección de la información.

\section{Estándares de seguridad}
\noindent
En general, el marco de referencia que se propone para este proyecto se basa en teorías y estándares reconocidos internacionalmente, ofreciendo así un soporte metodológico, técnico y estratégico para el diseño y desarrollo de la ampliación de la infraestructura HTCondor del grupo \GRID, permitiéndole a estos avanzar hacia una infraestructura más madura y escalable, alineada con las necesidades actuales de la computación distribuida y de alto rendimiento.
